\section{MA 与 ARMA：冲击的持久影响}
\label{sec:06-Random-Processes/06-ARMA-and-ARIMA/02}

AR 模型假设未来值由历史值直接加权决定。上一节你看到它的自相关按指数拖尾，偏自相关干净截尾，识别起来有章可循。但你在实际序列里会碰到另一种信号：一个外部冲击来了——政策调整、设备故障、市场恐慌——然后它的影响跨越多期逐渐退潮，但历史值本身并不直接出现在方程右边。比如一个工厂因为设备故障停产一天，次日出货量下降不是因为前一天出货少，而是因为同一次故障在三天的生产排期中都留下了痕迹。今天的问题想捕捉的就是这种``冲击的持久影响''。

AR 把影响储存在可观测的历史值里；MA 把影响储存在不可观测的历史冲击里。两者都不对，两者都有用——要看数据的记忆模式更接近哪一种。

\subsection{MA(q)：只用冲击说话}

MA(q) 模型直接把当前值写成历史白噪声的加权和：

\[
X_t=\varepsilon_t+\theta_1\varepsilon_{t-1}+\cdots+\theta_q\varepsilon_{t-q},
\]

其中 \(\varepsilon_t\) 是白噪声——均值为零、方差恒定、前后不相关。

与 AR 的差异很直观：AR 写 \(X_t=\phi X_{t-1}+\varepsilon_t\)，今天依赖昨天的观测值；MA 写 \(X_t=\varepsilon_t+\theta\varepsilon_{t-1}\)，今天只依赖今天和昨天的冲击，不引用任何可观测的系统状态。一个冲击在第 \(t\) 期进入系统，直接从 \(X_t\) 开始显现，然后通过系数 \(\theta_1\) 进入 \(X_{t+1}\)，再通过 \(\theta_2\) 进入 \(X_{t+2}\)，直到 \(q\) 期后完全消失——之前的 AR 模型里冲击的效应是按 \(\phi^k\) 指数衰减、永不完全归零的。

MA(q) 自动平稳：有限个白噪声的加权和，均值和方差永远是常数。但辨识的钥匙在自相关函数（ACF）上：

\[
\rho_k=0\quad\text{对 }k>q.
\]

ACF 在滞后 \(q\) 后截尾——这跟 AR 的 PACF 截尾形成了对偶关系：

\begin{itemize}
\tightlist
\item
  AR(p)：ACF 拖尾，PACF 在滞后 \(p\) 后截尾；
\item
  MA(q)：ACF 在滞后 \(q\) 后截尾，PACF 拖尾。
\end{itemize}

画出 ACF 和 PACF，前者在某个滞后之后突然归零、后者缓缓衰减，你手上的大概率是 MA 过程。反过来，PACF 截断、ACF 拖尾，那就是 AR。这套对偶是 Box-Jenkins 方法里阶数识别的核心直觉——不需要解方程，两张图就能先把候选模型圈出来。

\subsection{可逆性：MA 也有隐藏条件}

平稳性是 AR 的约束（特征根在单位圆外）；MA 天然平稳，但它有另一条关卡：可逆性。

MA(q) 写成滞后算子形式 \(X_t=\theta(B)\varepsilon_t\)，其中 \(\theta(B)=1+\theta_1B+\cdots+\theta_qB^q\)。可逆性要求 \(\theta(B)=0\) 的根在单位圆外。

满足可逆性时，MA 可以形式化地转为 AR(\(\infty\))——当前的冲击能用历史观测值线性表示：

\[
\varepsilon_t=\theta(B)^{-1}X_t.
\]

不满足会怎样？多个不同的参数化会产生完全相同的 ACF。给你一组数据，ACF 算出来一模一样，你分不出背后到底是哪组参数。这对参数估计是不稳定性的根源：似然面上有多个谷底。可逆性条件实际上是唯一性条件——它选出那个能用历史数据``反推''出来当前冲击的参数化。

\subsection{ARMA(p,q)：AR 和 MA 可以同居}

真实序列往往既有惯性也有滞后冲击的吸收。把两者拼在一起：

\[
X_t=\phi_1X_{t-1}+\cdots+\phi_pX_{t-p}
+\varepsilon_t+\theta_1\varepsilon_{t-1}+\cdots+\theta_q\varepsilon_{t-q}.
\]

用滞后算子简洁书写：\(\phi(B)X_t=\theta(B)\varepsilon_t\)。

平稳性由 \(\phi(B)\) 的根决定（跟纯 AR 一样）；可逆性由 \(\theta(B)\) 的根决定（跟纯 MA 一样）。两个条件各自独立检查，不需要联立求解。

辨识上的坏消息：ARMA 的 ACF 和 PACF 都拖尾，两者都不会干干净净地在某个滞后截断。单看这两张图只能排除纯 AR 或纯 MA，但给不出具体的 \(p\) 和 \(q\)。需要扩展 ACF（EACF）、信息准则（AIC/BIC），或者最原始的办法：先拟合几个候选模型，看残差的 ACF 像不像白噪声，再从候选中根据准则挑选。

以 ARMA(1,1) 为例感受一下混合效果：

\[
X_t=\phi X_{t-1}+\varepsilon_t+\theta\varepsilon_{t-1}.
\]

它的 ACF 为

\[
\rho_1=\frac{(1+\phi\theta)(\phi+\theta)}{1+\theta^2+2\phi\theta},
\qquad
\rho_k=\phi\rho_{k-1}\ (k\ge2).
\]

从滞后 2 开始按 \(\phi\) 指数衰减——这是 AR 的典型行为。但 \(\rho_1\) 由 \(\phi\) 和 \(\theta\) 共同决定，不单独属于 AR 那半或 MA 那半。``AR 决定衰减速率，MA 调整首个自相关值''——这是 ARMA 混合效应的经典截面。

\subsection{ARMA 与线性动力系统的深层联系}

ARMA 不是孤立的时间序列技巧。它和状态空间模型共享同一套数学骨架。

在状态空间表述中，系统有一个不可直接观测的状态向量，观测值是状态的线性组合加上噪声。ARMA(p,q) 可以精确地改写为一个状态空间模型——状态由历史值和历史冲击共同组成。反过来，任何线性高斯状态空间模型在稳态下都可以写成 ARMA 形式。

这个联系把 ARMA 的预测跟 Kalman 滤波对接上了：ARMA 的最小均方误差预测与 Kalman 对稳态系统的预测是等价的。状态空间框架的优势在于它天然处理时变参数、缺失数据和任意初始条件。ARMA 的优势在于它参数更少、识别更直观、理论更成熟。知道它们是同一枚硬币的两面，你就拥有了两套工具箱之间的便携转换手册。

\subsection{ARMA 画好的圈，和圈外的世界}

ARMA 假设系统是线性的、平稳的、创新是白噪声。以下任何一条不成立时，需要越过 ARMA 的边界：

\begin{itemize}
\tightlist
\item
  \textbf{非平稳}：存在单位根或确定性趋势，均值或方差随时间漂移——需要 ARIMA 或含趋势项的状态空间模型；
\item
  \textbf{长记忆}：ACF 不是指数衰减，而是按幂律极缓慢地拖尾——需要 ARFIMA（分数阶差分）；
\item
  \textbf{非线性}：条件均值对过去值的依赖不是线性的——需要 TAR、STAR 或丢掉参数化直接用神经网络；
\item
  \textbf{异方差}：创新的方差本身随时间变化、存在聚集效应——需要 GARCH 族；
\item
  \textbf{多元联动}：多个序列互相影响——需要 VAR（向量 AR）放开维数。
\end{itemize}

这些扩展各有自己的辨识逻辑和估计难点，但 ARMA 之所以仍然是起点，是因为它是``最简的完整模型''：ACF/PACF 给出干净直觉，Yule-Walker 和极大似然给出可靠估计，Box-Jenkins 方法给出系统化的建模流程。下一篇处理非平稳序列，就从最自然的扩展——允许差分——开始。
